\subsection{Semiconductor Data}
\label{sec:semiconductor}

The example module is the Mitsubishi Electric CM1000HG-130XA, 6.5\,kV, 1000\,A, with the data at
$T_j = 125\,^{\circ}$C~[2]. The on-state voltage of IGBT and diode is modelled per module as
\begin{equation}
  v_{CE} = v_{0\_igbt} + r_{igbt} \cdot i_C , \qquad v_{EC} = v_{0\_diode} + r_{diode} \cdot i_E ,
\end{equation}
fitted through the datasheet values at 1000\,A ($V_{CE\_sat} = 3.25$\,V, $V_{EC} = 2.80$\,V, typical, $T_j = 125\,^{\circ}$C)
and the typical output characteristics at 500\,A. This gives $v_{0\_igbt} = 1.45$\,V, $r_{igbt} = 1.80$\,m$\Omega$,
$v_{0\_diode} = 1.60$\,V and $r_{diode} = 1.20$\,m$\Omega$. Table~\ref{tab:onstate} compares the model with the
datasheet: the deviation is at most 0.1\,V from 250\,A to 1000\,A, which covers the module currents of the converter
(up to 0.69\,kA per module in the example), and reaches 0.5\,V at 2000\,A. The points other than 1000\,A are read
from the characteristics with an accuracy of about $\pm 0.1$\,V.

\begin{table}[!ht]
  \centering
  \small
  \begin{tabular}{@{}l r r r r@{}}
    \toprule
    \textbf{Device} & \multicolumn{1}{c}{\textbf{$i$ [kA]}} & \multicolumn{1}{c}{\textbf{Datasheet [V]}}
      & \multicolumn{1}{c}{\textbf{Model [V]}} & \multicolumn{1}{c}{\textbf{Deviation [V]}} \\
    \midrule
    \tablerows{thermal/fig/onstate_table}
    \bottomrule
  \end{tabular}
  \caption{On-state model of the CM1000HG-130XA against the datasheet, $T_j = 125\,^{\circ}$C~[2]}
  \label{tab:onstate}
\end{table}

The switching energies per pulse are those of the datasheet at $V_{CC} = 3600$\,V, $I_C = 1000$\,A,
$R_{G\_on} = 4.3\,\Omega$, $R_{G\_off} = 39\,\Omega$, $L_s = 150$\,nH and $T_j = 125\,^{\circ}$C, as defined in
IEC 60747: $E_{on} = 11.0$\,J, $E_{off} = 7.5$\,J, $E_{rec} = 5.1$\,J~[2]. The converter scales them linearly with
voltage and current (Section~\ref{sec:converter-losses}). The datasheet curves of the energies against current are
not linear; below 1000\,A the linear scaling gives lower energies than the curves.

The thermal data are $R_{th\_jc} = 10.0$\,K/kW (IGBT) and 16.0\,K/kW (diode), maximum values, the contact resistance
case to heat sink $R_{th\_cs} = 5.0$\,K/kW per module (typical, grease $\lambda = 1$\,W/(m$\cdot$K), 80\,$\mu$m), the
highest operating junction temperature $T_{jop} = 150\,^{\circ}$C, and the Foster model of the transient thermal
impedance junction to case, the same for IGBT and diode:
\begin{equation}
  Z_{th\_jc}(t) = R_{th\_jc} \cdot \sum_{i=1}^{4} r_i \cdot \left( 1 - \mathrm{e}^{-t/\tau_i} \right) ,
\end{equation}
with $r_i = 0.0096$, 0.1893, 0.4044, 0.3967 and $\tau_i = 0.1$\,ms, 5.8\,ms, 60.2\,ms, 351.2\,ms~[2].

With the ratings given, the converter warns when the highest submodule capacitor voltage reaches the blocking
voltage $V_{CES}$ and when the peak current per module $I_{arm\_peak} / n_{par}$ exceeds the rated collector current
$I_C$. With one module per switch position the example exceeds $I_C$ (1.38\,kA); with two it does not.
